% Preamble for Pandoc Generation
% The page style is the base; the book is layered on top.
%
% latex/page-style.sty       the page: base packages, geometry, running-head frame
% latex/theorem-frame.sty    the mdframed style and the \declaretheorem key the
%                            book's environments are framed with
% latex/latex-template.sty   the book's macros, colours and helpers
% latex/chapter-style.sty    chapters, \part banners, running heads
% latex/theorem-envs.sty     the environments filters/theorems.lua emits, in that
%                            frame, in the colours from the book's config

% ======================================================================================
% ONE DECISION PER SHARED PACKAGE
% ======================================================================================
% Sixteen packages are loaded by both the page style and the book. Most take no options
% and so cannot clash; these are the ones that do. Naming the options here, before
% anything pulls the packages in, means no load order can produce an option clash
% and there is a single place that says which settings win.

% xcolor: page-style loads it plain, latex-template wants dvipsnames. dvipsnames wins.
\PassOptionsToPackage{dvipsnames}{xcolor}
% hyperref: loaded exactly once, by templates/latex/template.tex, after
% latex/macros.tex (it reclaims \H, restored at the bottom of this file). These
% are the book's link settings, handed to that one load.
\PassOptionsToPackage{pdfusetitle,linktoc=page,allcolors=blue}{hyperref}
% mdframed: theorem-frame and latex-template both ask for framemethod=TikZ, so
% the second load is a no-op. Do NOT hand it the same option from here as well:
% mdframed reads the frame method by \input-ing a file, and it would do it twice.
% tcolorbox (most,breakable,skins,theorems) is loaded only by latex-template,
% and caption's skip=1em only by page-style, which is loaded first -- latex-template's
% later option-free load of caption is a subset and cannot clash.

% ======================================================================================
% THE BASE
% ======================================================================================
% Page geometry, running-head frame, paragraph style, tables and figures, and
% the packages they need. latex/page-style.sty is the one place the page is decided; it
% keeps the book's uniform 2.5cm margin, so the measure -- and every line break
% in 3,300 pages -- is unchanged.
\usepackage{latex/page-style}

\usepackage{bbm}

% ======================================================================================
% THE BOOK ON TOP
% ======================================================================================
% Macros, colours and helpers. page-style owns the page (skipgeometry), the pandoc
% template owns hyperref (skiphyperref), latex/chapter-style owns the headings
% (skiptitlesec) and latex/theorem-envs owns the environments (skipenv).
\usepackage[skipgeometry,skiphyperref,skiptitlesec,skipenv]{latex/latex-template}
\let\maketitle\relax

% List spacing (1.2x)
\usepackage{enumitem}
\setlist{itemsep=1.2em, parsep=0pt}

% Chapter styling (bookmode: two-level no box, three-level black box, correct font
% sizes) plus the \part banners and running heads, none of which page-style has.
\usepackage[bookmode]{latex/chapter-style}

% The environments filters/theorems.lua emits, in the book's framed style. Which
% environments exist, their names, counters and colours come from the book's
% config: the build writes them to _build/tmp/environments.tex, included right
% after this file. That covers the numbered forms, the unnumbered ones and the
% starred variants the preface uses (definition*, theorem*, ...).
\usepackage{latex/theorem-envs}

% Code listings (runnable cells) and function plots (components.lua)
\usepackage{listings}
\lstset{basicstyle=\ttfamily\small, columns=fullflexible, keepspaces=true, breaklines=true,
  frame=single, rulecolor=\color{black!25}, framesep=6pt, xleftmargin=6pt, xrightmargin=6pt,
  keywordstyle=\color{blue!60!black}\bfseries, commentstyle=\color{black!55}\itshape,
  stringstyle=\color{red!60!black}, showstringspaces=false, upquote=true}
\usepackage{pgfplots}
\pgfplotsset{compat=1.18}

% Unicode symbols typed directly in prose (e.g. the proof tags (⇒), (⊆))
\usepackage{newunicodechar}
\newunicodechar{⇒}{\ensuremath{\Rightarrow}}
\newunicodechar{⇐}{\ensuremath{\Leftarrow}}
\newunicodechar{⇔}{\ensuremath{\Leftrightarrow}}
\newunicodechar{→}{\ensuremath{\to}}
\newunicodechar{⊆}{\ensuremath{\subseteq}}
\newunicodechar{⊇}{\ensuremath{\supseteq}}
\newunicodechar{∈}{\ensuremath{\in}}
\newunicodechar{≤}{\ensuremath{\le}}
\newunicodechar{≥}{\ensuremath{\ge}}
\newunicodechar{≠}{\ensuremath{\ne}}
\newunicodechar{×}{\ensuremath{\times}}
\newunicodechar{①}{\textcircled{\scriptsize 1}}
\newunicodechar{②}{\textcircled{\scriptsize 2}}
\newunicodechar{③}{\textcircled{\scriptsize 3}}
\newunicodechar{④}{\textcircled{\scriptsize 4}}
\newunicodechar{⑤}{\textcircled{\scriptsize 5}}
\newunicodechar{⑥}{\textcircled{\scriptsize 6}}
\newunicodechar{⑦}{\textcircled{\scriptsize 7}}
\newunicodechar{⑧}{\textcircled{\scriptsize 8}}
\newunicodechar{⑨}{\textcircled{\scriptsize 9}}
\newunicodechar{∗}{\ensuremath{\ast}}
\newunicodechar{−}{\ensuremath{-}}
\newunicodechar{ᵀ}{\texorpdfstring{\ensuremath{^{\top}}}{T}}
% Greek letters and common symbols typed as Unicode in prose and headings (titles keep math out of
% \(...\), so a heading may say "Recursion on ω"); bookmarks keep the character itself
\newunicodechar{α}{\texorpdfstring{\ensuremath{\alpha}}{α}}
\newunicodechar{β}{\texorpdfstring{\ensuremath{\beta}}{β}}
\newunicodechar{γ}{\texorpdfstring{\ensuremath{\gamma}}{γ}}
\newunicodechar{δ}{\texorpdfstring{\ensuremath{\delta}}{δ}}
\newunicodechar{ε}{\texorpdfstring{\ensuremath{\varepsilon}}{ε}}
\newunicodechar{ζ}{\texorpdfstring{\ensuremath{\zeta}}{ζ}}
\newunicodechar{η}{\texorpdfstring{\ensuremath{\eta}}{η}}
\newunicodechar{θ}{\texorpdfstring{\ensuremath{\theta}}{θ}}
\newunicodechar{ι}{\texorpdfstring{\ensuremath{\iota}}{ι}}
\newunicodechar{κ}{\texorpdfstring{\ensuremath{\kappa}}{κ}}
\newunicodechar{λ}{\texorpdfstring{\ensuremath{\lambda}}{λ}}
\newunicodechar{μ}{\texorpdfstring{\ensuremath{\mu}}{μ}}
\newunicodechar{ν}{\texorpdfstring{\ensuremath{\nu}}{ν}}
\newunicodechar{ξ}{\texorpdfstring{\ensuremath{\xi}}{ξ}}
\newunicodechar{π}{\texorpdfstring{\ensuremath{\pi}}{π}}
\newunicodechar{ρ}{\texorpdfstring{\ensuremath{\rho}}{ρ}}
\newunicodechar{σ}{\texorpdfstring{\ensuremath{\sigma}}{σ}}
\newunicodechar{τ}{\texorpdfstring{\ensuremath{\tau}}{τ}}
\newunicodechar{υ}{\texorpdfstring{\ensuremath{\upsilon}}{υ}}
\newunicodechar{φ}{\texorpdfstring{\ensuremath{\varphi}}{φ}}
\newunicodechar{χ}{\texorpdfstring{\ensuremath{\chi}}{χ}}
\newunicodechar{ψ}{\texorpdfstring{\ensuremath{\psi}}{ψ}}
\newunicodechar{ω}{\texorpdfstring{\ensuremath{\omega}}{ω}}
\newunicodechar{Γ}{\texorpdfstring{\ensuremath{\Gamma}}{Γ}}
\newunicodechar{Δ}{\texorpdfstring{\ensuremath{\Delta}}{Δ}}
\newunicodechar{Θ}{\texorpdfstring{\ensuremath{\Theta}}{Θ}}
\newunicodechar{Λ}{\texorpdfstring{\ensuremath{\Lambda}}{Λ}}
\newunicodechar{Ξ}{\texorpdfstring{\ensuremath{\Xi}}{Ξ}}
\newunicodechar{Π}{\texorpdfstring{\ensuremath{\Pi}}{Π}}
\newunicodechar{Σ}{\texorpdfstring{\ensuremath{\Sigma}}{Σ}}
\newunicodechar{Φ}{\texorpdfstring{\ensuremath{\Phi}}{Φ}}
\newunicodechar{Ψ}{\texorpdfstring{\ensuremath{\Psi}}{Ψ}}
\newunicodechar{Ω}{\texorpdfstring{\ensuremath{\Omega}}{Ω}}
\newunicodechar{′}{\texorpdfstring{\ensuremath{{}^{\prime}}}{′}}
\newunicodechar{″}{\texorpdfstring{\ensuremath{{}^{\prime\prime}}}{″}}
\newunicodechar{ℕ}{\texorpdfstring{\ensuremath{\mathbb{N}}}{ℕ}}
\newunicodechar{ℤ}{\texorpdfstring{\ensuremath{\mathbb{Z}}}{ℤ}}
\newunicodechar{ℚ}{\texorpdfstring{\ensuremath{\mathbb{Q}}}{ℚ}}
\newunicodechar{ℝ}{\texorpdfstring{\ensuremath{\mathbb{R}}}{ℝ}}
\newunicodechar{ℂ}{\texorpdfstring{\ensuremath{\mathbb{C}}}{ℂ}}
\newunicodechar{∀}{\texorpdfstring{\ensuremath{\forall}}{∀}}
\newunicodechar{∃}{\texorpdfstring{\ensuremath{\exists}}{∃}}
\newunicodechar{¬}{\texorpdfstring{\ensuremath{\neg}}{¬}}
\newunicodechar{∧}{\texorpdfstring{\ensuremath{\wedge}}{∧}}
\newunicodechar{∨}{\texorpdfstring{\ensuremath{\vee}}{∨}}
\newunicodechar{⊥}{\texorpdfstring{\ensuremath{\bot}}{⊥}}
\newunicodechar{⊢}{\texorpdfstring{\ensuremath{\vdash}}{⊢}}
\newunicodechar{⊨}{\texorpdfstring{\ensuremath{\models}}{⊨}}
\newunicodechar{∉}{\texorpdfstring{\ensuremath{\notin}}{∉}}
\newunicodechar{∅}{\texorpdfstring{\ensuremath{\emptyset}}{∅}}
\newunicodechar{∞}{\texorpdfstring{\ensuremath{\infty}}{∞}}
\newunicodechar{±}{\texorpdfstring{\ensuremath{\pm}}{±}}
\newunicodechar{⊂}{\texorpdfstring{\ensuremath{\subset}}{⊂}}
\newunicodechar{∪}{\texorpdfstring{\ensuremath{\cup}}{∪}}
\newunicodechar{∩}{\texorpdfstring{\ensuremath{\cap}}{∩}}
\newunicodechar{↔}{\texorpdfstring{\ensuremath{\leftrightarrow}}{↔}}
\newunicodechar{↦}{\texorpdfstring{\ensuremath{\mapsto}}{↦}}
\newunicodechar{∘}{\texorpdfstring{\ensuremath{\circ}}{∘}}
\newunicodechar{√}{\texorpdfstring{\ensuremath{\surd}}{√}}
\newunicodechar{≈}{\texorpdfstring{\ensuremath{\approx}}{≈}}
\newunicodechar{≡}{\texorpdfstring{\ensuremath{\equiv}}{≡}}
\newunicodechar{⋯}{\texorpdfstring{\ensuremath{\cdots}}{⋯}}
\newunicodechar{∣}{\texorpdfstring{\ensuremath{\mid}}{∣}}

\newunicodechar{ }{\,}  % U+2009 thin space, as between initials ("E. H. Moore")

% Images written by Pandoc carry both width and height=\textheight; keep their shape
\usepackage{graphicx}
\setkeys{Gin}{keepaspectratio}

% Tables written by Pandoc (column widths use \real from calc)
\usepackage{longtable,booktabs,array,calc}
\providecommand{\tightlist}{\setlength{\itemsep}{0pt}\setlength{\parskip}{0pt}}

% hyperref is loaded by the pandoc template after latex/macros.tex, and it
% reclaims \H (the Hungarian-umlaut accent). Restore the bold matrix H once
% the preamble is over, so \H means the bold matrix H -- in math. In the prose
% \H is still the accent, or "Szeg\H{o}" loses its o (\bt@textsafe is
% latex/latex-template.sty's; this hook is the last one to touch \H).
\makeatletter
\AtBeginDocument{\bt@textsafe\H\mH}
\makeatother

% Let TeX stretch a line rather than run it into the margin. Measured on two books
% (linear algebra, regression) over three full builds each: 1em removes 54-61% of the
% body's overfull boxes and changes no label number and no page count. 3em goes further
% but re-breaks subsection headings onto conspicuously loose lines, which reads worse
% than the overrun it fixes. The Contents' overfull boxes are a separate cause; see
% latex/chapter-style.sty (\cftsubsecnumwidth).
\emergencystretch=1em

% Fitch-style derivations, if the book ships latex/fitch-book.sty (inert otherwise;
% filters/tikz.lua and build/tikz.py handle the `nd` environment it defines).
\IfFileExists{fitch-book.sty}{\usepackage{fitch-book}}{}
